% ============================================================
\chapter{Literature Survey}
\label{chap:lit}
% ============================================================

This chapter surveys work in the areas PipedPeer draws on: frameworks for
distributed Python and low-communication training, lightweight sandboxing and
environment portability at the edge, actor-style fault tolerance,
failure-transparent programming, portability in the Python ecosystem,
scheduling, and collaborative task offloading. Each review closes with its
bearing on the design of Chapters~\ref{chap:arch} and~\ref{chap:design}.

``Spark: Cluster Computing with Working Sets'' \cite{zaharia2010spark} by
M.~Zaharia \textit{et al.} keeps a working set in memory across parallel
operations as resilient distributed datasets, which a lost partition rebuilds
from its lineage, and runs iterative jobs up to ten times faster than
MapReduce \cite{dean2004mapreduce}. A program is written against the dataset
API and runs on a cluster with a master, which is exactly the set-up PipedPeer
removes; its hash partitioning by key is the technique PipedPeer's dataframe
shuffle uses (Section~\ref{sec:intercept}).

``Ray: A Distributed Framework for Emerging AI Applications''
\cite{moritz2018ray} by P.~Moritz \textit{et al.} unifies tasks and actors over
a distributed scheduler and a sharded control store, and reaches millions of
tasks per second. Ray offers a drop-in \texttt{multiprocessing.Pool} and a
\texttt{joblib} back-end, the closest existing answer to PipedPeer's code
problem, but each needs a changed import and a running cluster with a head
node on one network. PipedPeer substitutes the same primitives without the
import and decides per call whether distributing them pays.

``Dask: Parallel Computation with Blocked Algorithms and Task Scheduling''
\cite{rocklin2015dask} by M.~Rocklin expresses NumPy- and pandas-like
operations as graphs of blocked tasks run by a dynamic scheduler, so that
arrays and dataframes larger than memory can be processed in parallel. Its
interfaces mirror the libraries they replace, yet the program must still be
written against Dask's collections. PipedPeer's out-of-core \texttt{read\_csv}
and row-block matrix products follow Dask's blocking idea while leaving the
original pandas and NumPy calls in place.

``DiLoCo: Distributed Low-Communication Training of Language Models''
\cite{douillard2023diloco} by A.~Douillard \textit{et al.}, the work closest to
PipedPeer's training path, trains replicas independently for hundreds of steps
and merges their parameter changes with an outer Nesterov step, matching
synchronous training on a language model while communicating about 500 times
less. PipedPeer's fast synchronisation takes the same structure for links
between homes, but its intervals are 2 to 64 steps rather than hundreds, and at
those intervals the outer momentum multiplies the script's learning rate; it
therefore averages plainly (Section~\ref{sec:ddp}). BOINC
\cite{anderson2004boinc}, finally, showed that volunteers' idle machines can
serve real computation, but as a platform for specific projects with their own
servers rather than a general executor.

``Web Application-Based WebAssembly Container Platform for Extreme Edge
Computing'' \cite{sekigawa2023} by S.~Sekigawa \textit{et al.} deploys
WebAssembly modules as isolated compute units on edge nodes, sandboxing and
executing user-submitted binaries without a container daemon, and shows
acceptable latency with strong isolation on constrained hardware. Its argument
that a lightweight sandbox can replace a full container runtime is the one
PipedPeer follows, but PipedPeer must run arbitrary CPython with native
extensions, which WebAssembly cannot host, so it uses an OCI runtime
(\texttt{crun}) directly; Section~\ref{sec:res-sandbox} measures the cost of
that choice at about 20\,ms per job.

``WebAssembly for Edge Computing: Potential and Challenges'' \cite{hoque2022}
by M.~N.~Hoque and K.~A.~Harras weighs WebAssembly's near-native speed and memory
sandboxing against containers and virtual machines, and identifies an immature
system interface, limited file-system and networking APIs, and absent GPU support
as the blocking gaps. These gaps are why PipedPeer ships a complete native
environment as a Nix closure rather than compiling user code to a portable
bytecode: a closure carries native dependencies and GPU libraries intact.

``A Cross-Architecture Evaluation of WebAssembly in the Cloud-Edge Continuum''
\cite{kakati2024} by S.~Kakati and M.~Brorsson benchmarks WebAssembly runtimes
across x86\_64, ARM64 and RISC-V, reporting 70--90\% of native performance and
start-up under 10\,ms, with I/O-bound workloads suffering most from system-call
translation. PipedPeer also targets heterogeneous architectures but resolves
them at build time: the flake generator emits an architecture-specific closure,
and architecture is a hard filter in placement (Section~\ref{sec:placement}), so
full native performance is retained.

``Design of Elixir-Based Edge Server for Responsive IoT Applications''
\cite{li2023elixir} by Y.~Li and S.~Fujita builds an IoT edge server on Elixir and
OTP, where supervisor trees restart failed processes and the system stays
available through individual crashes. PipedPeer's retry loop follows the same
model, placing a failed job again on the next-best node with exponential
back-off, but it has no supervisor process: recovery follows from lease expiry
(Section~\ref{sec:lease}), so there is nothing whose own failure needs
supervising.

``A Synergistic Elixir-EDA-MQTT Framework for Smart Transportation''
\cite{zhou2024} by the same authors combines Elixir concurrency, an event-driven
architecture and MQTT publish-subscribe messaging, showing lower overhead than
REST for high-frequency telemetry. PipedPeer's traffic is of a different kind:
closures, workspaces and results are bulk transfers that need streaming and
back-pressure, so dispatch uses plain HTTP with a WebSocket for output rather
than a message broker.

``WebAssembly with wasi-nn for Edge Machine Learning Inference''
\cite{bachmeier2025} by J.~Bachmeier \textit{et al.} finds good portability for CPU
inference under WebAssembly and immaturity for GPU acceleration. Machine
learning is a primary use case for PipedPeer, which is why torch carries the
largest entry in its memory estimator and why GPU passthrough is implemented at
the OCI level; the paper's negative result on GPU support supports shipping a
native CUDA-capable environment instead.

``Case Study of WebAssembly Runtimes for AI Applications on the Edge''
\cite{khelifa2024} by S.~E.~Khelifa \textit{et al.} compares four runtimes for AI
inference on edge devices and finds memory consumption to be the dominant
constraint. The same assumption underlies PipedPeer's memory estimator and
admission control: a node refuses a job it cannot hold, a parallel-map request
that would exceed 40\% of a node's memory is split into sequential parts, and
every job runs under its own memory ceiling inside a daemon that is itself
capped (Section~\ref{sec:sandbox}).

``Programming Models for Failure-Transparent Distributed Systems''
\cite{kth2025thesis} by J.~Spenger defines failure transparency --- masking
failures fully from the programmer --- independently of a particular semantics,
proves it for a model of stateful dataflow streaming, and builds two programming
models on that base. Failure transparency is PipedPeer's central invariant: the
user's script contains no failure handling, so every remote path must degrade to
a working local one. PipedPeer enforces this by layered fallback and tests rather
than by proof (Sections~\ref{sec:ft} and~\ref{sec:res-ft}).

``Python Array API Standard'' \cite{arrayapi2023} by A.~Meurer \textit{et al.}
standardises an array interface so that the same code runs on NumPy, CuPy or JAX
without modification. Its goal, portability without code change, is
PipedPeer's, but the mechanism is the mirror image: the standard asks libraries
to converge on one interface that users then adopt, while PipedPeer changes
nothing the user writes and substitutes the implementations behind the
interfaces the script already uses.

``iDDS: Intelligent Distributed Dispatch and Scheduling'' \cite{idds2025} by
W.~Guan \textit{et al.} assigns workflow tasks to heterogeneous resources from
real-time CPU, memory, data-locality and dependency constraints, driven by a
task graph. Its scoring-based assignment is comparable to PipedPeer's placement
function (Section~\ref{sec:placement}), with two differences: PipedPeer schedules
whole scripts rather than task graphs, and its scoring runs in the submitting
client, with no scheduler process. Data locality, which iDDS treats as
first-class, is future work for PipedPeer.

``Task Offloading in Multi-Hop Relay-Aided Multi-Access Edge Computing''
\cite{deng2023} by Y.~Deng \textit{et al.} formulates joint task assignment and
relay selection to minimise completion time under transmission and energy
constraints. Its comparison of transmission time against local computation time
has the same shape as PipedPeer's cost model (Section~\ref{sec:cost}); PipedPeer
measures link bandwidth at run time rather than assuming a channel model, and for
parallel maps fits each peer's set-up time and throughput from the work it has
returned.

``Multi-Agent Deep Reinforcement Learning Based Dynamic Task Offloading''
\cite{he2024} by H.~He \textit{et al.} lets each device learn an offloading policy
from local observations without a central coordinator, outperforming centralised
heuristics on the energy-delay trade-off. A learned policy needs training data
and a training phase, which is at odds with a tool that a new machine joins by
running one binary; PipedPeer therefore uses a deterministic rule-based model,
and a learned policy is a natural direction once a deployment has history.

``Joint Optimization of Multi-User Partial Offloading and Resource Allocation''
\cite{yong2023} by D.~Yong \textit{et al.} optimises the split between locally
executed and offloaded portions of a task and finds partial offloading
consistently better than all-or-nothing. PipedPeer offloads partially in this
sense: interception splits a script's work at the granularity of pool chunks,
shuffle buckets and matrix row blocks while the rest runs locally, and the split
is chosen by minimising completion time over measured rates
(Section~\ref{sec:intercept}).

``Energy-Efficient Collaborative Task Offloading in MEC'' \cite{adhoc2025} by
S.~Wang \textit{et al.} has devices cooperate by relaying tasks toward edge
servers, with a learned policy that cuts total energy. Its symmetric model, in
which every device is both client and server, is PipedPeer's: one binary, and any
node may submit or accept work. PipedPeer likewise forwards a parallel-map
request that a node cannot hold to a peer that can.

``Blockchain Enabled Task Offloading in Digital Twin Vehicular Edge Networks''
\cite{li2023blockchain} by C.~Li \textit{et al.} enforces service agreements between
task submitters and edge workers with smart contracts, preventing free-riding in
an untrusted network. PipedPeer authenticates who may use a node, with a shared
secret, encrypted links and signed closures (Section~\ref{sec:trust}), but does
not address whether an admitted participant behaves well; the paper describes
the class of mechanism a cluster of mutually untrusting owners would need.

``Congestion-Aware Distributed Task Offloading Using Graph Neural Networks''
\cite{zhao2023} by Z.~Zhao \textit{et al.} trains a graph neural network to produce
distributed offloading policies from local neighbourhood observations. PipedPeer
scores each node on scalar metrics and is blind to topology when it places a
whole job; link cost is measured only where work is divided. Feeding a
graph-structured view into placement is the most concrete extension this survey
suggests.

The survey leaves a gap that PipedPeer fills. The systems for distributed Python
--- Spark, Ray, Dask --- assume a configured cluster on one network; the
offloading literature models the decision of what to move but not the
environment, the program or the network path; and the sandboxing work stops at a
single node. None runs an unmodified script across machines behind different
routers.

\textbf{Standards for cross-network operation.} Joining machines on different
networks (Section~\ref{sec:nat}) rests on established standards. Address
discovery uses the reflexive-address idea of STUN \cite{rfc8489}, and a router's
mapping and filtering behaviour is classified in the terms of RFC~5780
\cite{rfc5780}. Port mapping follows NAT-PMP \cite{rfc6886}, its successor PCP
\cite{rfc6887}, and UPnP IGD \cite{upnpigd}. Simultaneous UDP probing from the
socket that will carry the connection is the hole-punching technique analysed by
Ford, Srisuresh and Kegel \cite{ford2005p2p}, whose observation that a router's
mapping belongs to a socket shapes PipedPeer's shared-socket design. Links are
QUIC \cite{rfc9000} secured by TLS~1.3 \cite{rfc8446}, node identities are
Ed25519 keys \cite{bernstein2012ed25519}, and archives are compressed with
Zstandard \cite{rfc8878}.
